\section{Tesseract-BP Architecture \& Optimization Deep-Dive}

\begin{frame}{Tesseract-BP High-Level Data Flow}
    \centering
    \begin{tikzpicture}[node distance=1.2cm, auto, >=stealth]
        \node [rectangle, draw, fill=googleblue!15, rounded corners, text width=2.5cm, align=center] (dem) {\textbf{Stim DEM}\\\scriptsize Circuit Error Model};
        \node [rectangle, draw, fill=googleblue!30, rounded corners, text width=2.5cm, align=center, right=0.6cm of dem] (csr) {\textbf{CSR Tanner Graph}\\\scriptsize Sparse Checks/Vars};
        \node [rectangle, draw, fill=googleyellow!25, rounded corners, text width=2.8cm, align=center, right=0.6cm of csr] (bp) {\textbf{AVX-512 BP Engine}\\\scriptsize Layered Serial Min-Sum\\(Batches of 16)};
        \node [diamond, draw, fill=googleyellow!50, align=center, right=0.6cm of bp, aspect=1.5] (check) {\scriptsize Converged?};
        \node [rectangle, draw, fill=googlegreen!30, rounded corners, text width=2.2cm, align=center, above right=0.4cm and 0.8cm of check] (out) {\textbf{Hard Decision}\\\scriptsize Syndrome Cleared};
        \node [rectangle, draw, fill=googlered!25, rounded corners, text width=2.5cm, align=center, below right=0.4cm and 0.8cm of check] (osd) {\textbf{OSD Solver}\\\scriptsize Gaussian Elimination};
        
        \draw [->, thick] (dem) -- (csr);
        \draw [->, thick] (csr) -- (bp);
        \draw [->, thick] (bp) -- (check);
        \draw [->, thick] (check) -- node[above] {\scriptsize Yes} (out);
        \draw [->, thick] (check) -- node[below] {\scriptsize No} (osd);
        \draw [->, thick] (osd) |- (out);
    \end{tikzpicture}
\end{frame}

\begin{frame}[fragile]{Architecture: Pointer-Chasing vs. Bidirectional CSR}
    \begin{columns}[T]
        \begin{column}{0.5\textwidth}
            \begin{block}{\textcolor{googlered}{Naive BP}: Pointer-Chasing Graph}
                \tiny
                \begin{lstlisting}[language=C++]
struct Node {
    std::vector<Node*> neighbors;
    std::vector<double> messages;
};
std::vector<Node*> nodes;
                \end{lstlisting}
                \normalsize
                \begin{itemize}
                    \item Follows \texttt{std::vector} memory allocations.
                    \item Nodes scattered across fragmented heap memory.
                    \item Traversing neighbors means chasing pointers $\rightarrow$ \textbf{L1/L2 Cache Miss Penalty}!
                \end{itemize}
            \end{block}
        \end{column}
        
        \begin{column}{0.5\textwidth}
            \begin{block}{\textcolor{googlegreen}{Tesseract BP}: Bidirectional CSR}
                \tiny
                \begin{lstlisting}[language=C++]
// Flat integer arrays, no pointers!
vector<int> check_to_var_indices;
vector<int> check_to_var_offsets;
vector<int> var_to_check_indices;
vector<int> var_to_check_offsets;
                \end{lstlisting}
                \normalsize
                \begin{itemize}
                    \item \textbf{Compressed Sparse Row (CSR)}: Graph represented as 1D contiguous arrays.
                    \item \textbf{Bidirectional}: Fast $O(1)$ neighborhood lookups in both directions ($C \rightarrow V$ and $V \rightarrow C$).
                    \item Sequential memory access ensures hardware prefetcher efficiency.
                \end{itemize}
            \end{block}
        \end{column}
    \end{columns}
\end{frame}

\begin{frame}[fragile]{Optimization 1: Memory Foundation Rewrite (Interleaved 1D Array)}
    \begin{columns}[T]
        \begin{column}{0.5\textwidth}
            \begin{block}{\textcolor{googlered}{BEFORE}: Vector of Vectors (Non-Interleaved)}
                \tiny
                \begin{verbatim}
Shot 0:  [ Var0, Var1, Var2, ..., VarN ]
Shot 1:  [ Var0, Var1, Var2, ..., VarN ]  (RAM Gap)
Shot 2:  [ Var0, Var1, Var2, ..., VarN ]  (RAM Gap)
...
Shot 15: [ Var0, Var1, Var2, ..., VarN ]
                \end{verbatim}
                \normalsize
                \begin{itemize}
                    \item Loading Variable $V_i$ across 16 SIMD shots required fetching \textbf{16 separate RAM addresses}.
                    \item Transfers 1,024 bytes to extract 64 bytes of data $\rightarrow$ \textbf{L1 Cache Flushed!}
                \end{itemize}
            \end{block}
        \end{column}
        
        \begin{column}{0.5\textwidth}
            \begin{block}{\textcolor{googlegreen}{AFTER}: Flat Interleaved 1D Array}
                \tiny
                \begin{verbatim}
RAM: [ Var0_S0..S15 | Var1_S0..S15 | Var2_S0..S15 ]
                \end{verbatim}
                \normalsize
                \begin{itemize}
                    \item Memory address: \texttt{posteriors\_flat[v * 16 + b]}.
                    \item All 16 shot streams for Variable $V_i$ sit \textbf{physically contiguous} in RAM (64 bytes).
                    \item \textbf{1 SIMD Aligned Load} = 1 L1 Cache Line Hit!
                \end{itemize}
            \end{block}
        \end{column}
    \end{columns}
\end{frame}

\begin{frame}{Optimization 2: Horizontal Layered Serial Scheduling}
    \begin{block}{Algorithmic Transformation}
        Instead of updating variable node posteriors sequentially, iteration loops operate directly over \textbf{check nodes ($c \in C$)}:
        \begin{enumerate}
            \item For each check node $c$, calculate check-to-variable updates dynamically:
                  $$Q_{c,v} = L_v - R_{c,v}$$
            \item Compute the updated check message $R'_{c,v}$ using Min-Sum across check neighbors.
            \item Immediately update the variable posteriors in RAM:
                  $$L'_v = Q_{c,v} + R'_{c,v}$$
        \end{enumerate}
    \end{block}
    
    \begin{block}{50\% Memory Bandwidth Reduction}
        \begin{itemize}
            \item By performing $L_v - R_{c,v}$ on the fly, the explicit 2D message matrix $R_{c,v}$ is \textbf{completely eliminated from memory}!
            \item Reduces BP memory payload by \textbf{50\%}, doubling memory bus efficiency.
        \end{itemize}
    \end{block}
\end{frame}

\begin{frame}{Optimization 3: Explicit AVX-512 SIMD Vectorization}
    \begin{block}{Vector Register Mapping}
        16 parallel decoding shots ($32$-bit integer LLRs) map directly into one 512-bit AVX-512 register (\texttt{\_\_m512i}):
        $$16 \times 32 \text{ bits} = 512 \text{ bits}$$
    \end{block}
    
    \begin{block}{Core AVX-512 Intrinsic Operations}
        \begin{itemize}
            \item \texttt{\_mm512\_abs\_epi32}: Compute absolute LLR magnitudes for 16 shots simultaneously.
            \item \texttt{\_mm512\_min\_epi32}: Track minimum LLRs (\texttt{min1}, \texttt{min2}) without conditional branching latency.
            \item \texttt{\_mm512\_sub\_epi32} / \texttt{\_mm512\_add\_epi32}: Vectorized posterior updates.
            \item \texttt{\_mm512\_mask\_blend\_epi32} / \texttt{\_mm512\_movepi32\_mask}: Masked updates for active unconverged shots.
        \end{itemize}
    \end{block}
\end{frame}

\begin{frame}[fragile]{Optimization 4: Stochastic Check-Node Schedule Permutations}
    \begin{columns}[T]
        \begin{column}{0.55\textwidth}
            \begin{block}{Breaking Trapping Sets with PRNG}
                \begin{itemize}
                    \item Fixed schedule $0 \rightarrow M-1$ locks into symmetric LDPC trapping sets.
                    \item \textbf{Solution}: Apply Fisher-Yates check-node permutation before each BP iteration using a fast \textbf{Xorshift64 PRNG}.
                    \item Because memory is interleaved in 64-byte blocks, random check node hopping incurs \textbf{zero SIMD alignment penalty}!
                    \item Resolves trapping sets in Color Codes and Bivariate Bicycle codes seamlessly.
                \end{itemize}
            \end{block}
        \end{column}
        
        \begin{column}{0.42\textwidth}
            \begin{block}{Xorshift64 PRNG}
                \scriptsize
                \begin{lstlisting}[language=C++]
inline uint64_t xorshift64(uint64_t &state) {
    uint64_t x = state;
    x ^= x << 13;
    x ^= x >> 7;
    x ^= x << 17;
    return state = x;
}
                \end{lstlisting}
            \end{block}
        \end{column}
    \end{columns}
\end{frame}

\begin{frame}{Optimization 5: Ordered Statistic Decoding (OSD) Integration}
    \begin{block}{Algebraic Fallback Pipeline}
        When BP fails to clear all syndrome bits after \texttt{max\_iter}:
        \begin{enumerate}
            \item \textbf{Reliability Ordering}: Sort variable nodes by absolute LLR magnitude $|L_v|$ (highest magnitude = highest reliability).
            \item \textbf{Information Set Basis}: Perform column-pivoted Gaussian elimination on $H$ matrix over SIMD bitsets (\texttt{stim::simd\_bits}).
            \item \textbf{Perturbation Search}:
                  \begin{itemize}
                      \item \textbf{OSD-0}: Hard decision on basis + solver solution.
                      \item \textbf{OSD-1}: Test 1-bit perturbations on least reliable columns.
                      \item \textbf{OSD-2}: Test 2-bit perturbations on candidate combinations.
                  \end{itemize}
            \item Select candidate matching syndrome $s$ with minimum total weight/cost.
        \end{enumerate}
    \end{block}
\end{frame}
